\chapter{Reading Market Data}\label{ch:m1:reading-market-data}

Before lunch an exchange has sent one firm tens of millions of messages about
thousands of instruments. Each says that an order was added, reduced,
executed or removed, in a few dozen bytes. None says what the best bid is:
that is for the receiver to work out, by applying every message, in order, to
a book it maintains itself. If one message is lost, the book is wrong until
the end of the day, silently. If one trade report is a typing error, a price
of 10.01 in a stock worth 100, the day's low, its volatility, its
volume-weighted average price and every signal computed from them are wrong,
also silently. Every result in the rest of this series stands on a table of
prices; this chapter is about how that table is made, and about the defects
it inherits.

\section{Three levels of data}

\begin{definition}[Level 1, level 2, level 3]\label{def:m1:reading-market-data:levels}
\emph{Level 1}\index{level 1} data is the best bid and offer with their
sizes, and the last trade. \emph{Level 2}\index{level 2} data is the book
aggregated by price: for each price level, the total size (and sometimes
the number of orders), to some depth or in full: the \omterm{def:m1:matching-algorithms-and-implied-spreads:mbo}{market-by-price} feed of
\cref{def:m1:matching-algorithms-and-implied-spreads:mbo}. \emph{Level
3}\index{level 3} data is every order individually, with its identifier and
every event in its life: the \omterm{def:m1:matching-algorithms-and-implied-spreads:mbo}{market-by-order} feed, from which the other two
can be computed and which alone gives queue positions.
\end{definition}

\begin{omfigure}
\begin{tikzpicture}[font=\small,
  col/.style={draw=omDef,rounded corners=2pt,minimum width=3.9cm,minimum height=3.3cm,align=left,fill=omDef!4,font=\footnotesize\ttfamily,inner sep=5pt}]
\node[col] (l3) at (0,0) {ref side  size  price\\1001 B \ \ \ 300 \ 99.99\\1002 B \ \ \ 200 \ 99.99\\1004 B \ \ \ 100 \ 99.98\\1003 S \ \ \ 500 100.01\\1005 S \ \ \ 100 100.02};
\node[col] (l2) at (4.6,0) {side  size  price  n\\B \ \ \ \ 500 \ 99.99 \ 2\\B \ \ \ \ 100 \ 99.98 \ 1\\S \ \ \ \ 500 100.01 \ 1\\S \ \ \ \ 100 100.02 \ 1};
\node[col] (l1) at (9.2,0) {bid \ 99.99 x 500\\ask 100.01 x 500\\last 100.01 x 100};
\node[font=\small\bfseries,above=2pt of l3] {Level 3: orders};
\node[font=\small\bfseries,above=2pt of l2] {Level 2: prices};
\node[font=\small\bfseries,above=2pt of l1] {Level 1: top};
\draw[-{Stealth[length=2mm]},thick] (l3.east) -- (l2.west);
\draw[-{Stealth[length=2mm]},thick] (l2.east) -- (l1.west);
\end{tikzpicture}

\omcaption{One book at three levels of detail. Each level can be computed
from the one to its left and never the reverse: a vendor's one-second bars
are \omterm{def:m1:reading-market-data:levels}{level 1} sampled; they cannot be turned back into a
book.}\label{fig:m1:reading-market-data:levels}
\end{omfigure}

\begin{definition}[Sequence number]\label{def:m1:reading-market-data:seq}
A \emph{sequence number}\index{sequence number} is a counter that the
publisher increments with each message on a channel. A receiver that sees a
jump has lost messages and must recover them, from a retransmission service
or a snapshot, before its book can be trusted again.
\end{definition}

A \omterm{def:m1:reading-market-data:levels}{level 3} feed is a stream of \emph{differences}. The state is on the
receiver's side, which makes the protocol compact and makes every receiver
responsible for correctness. \omterm{def:m1:us-equity-market-structure:sip}{Direct feeds} are published on two redundant
network paths; the receiver takes whichever copy of each \omterm{def:m1:reading-market-data:seq}{sequence number}
arrives first and notices when neither does.

\begin{dated}{2026-09}{One wire format}\label{dat:m1:reading-market-data:itch}
The order-by-order equity feed whose version 5.0 specification is used as
this chapter's model encodes all integers big-endian; prices as integers
with four implied decimal places; timestamps as nanoseconds since midnight
in six bytes; and instruments by a two-byte \emph{locate} code assigned for
the day, at the same offset in every message. An add-order message has 36
bytes: type, locate, tracking number, timestamp, an eight-byte order
reference, side, shares, an eight-character symbol and the price. An
execution has 31 bytes and refers to the order only by its reference, as do
the cancel (23 bytes) and the delete (19 bytes): the price of a trade is not
in the execution message, it is wherever the receiver stored the order.
\end{dated}

\section{Trades, quotes and condition codes}

\begin{definition}[Trade condition]\label{def:m1:reading-market-data:condition}
A \emph{trade condition}\index{trade condition} (or sale condition) is a code
attached to a trade report that says what kind of trade it was and how it
may be used: a regular trade; an opening or closing auction print; an \omterm{def:m1:us-equity-market-structure:lots}{odd
lot}; a trade reported late or out of sequence; a trade at an average or
derived price; a correction or a cancellation of an earlier report. The
consolidated tapes and each exchange define which conditions update the last
price, the high and low, and the volume.
\end{definition}

A tape is not a list of prices at which one could have traded. An
out-of-sequence report carries a price from minutes ago. An average-price
trade matches no quote that ever existed. A cancelled trade did not happen;
its cancellation arrives as a later message that a careless loader ignores.
An auction print is one price for a large share of the day's volume. Off-exchange
trades (\cref{ch:m1:retail-flow-and-wholesaling}) are reported after the
fact, seconds late, with a timestamp that may be the reporter's.
Before any statistic is computed from trades, a rule must say which
conditions are in.

\begin{method}[Cleaning a tape]\label{met:m1:reading-market-data:clean}
\begin{enumerate}
\item Apply corrections and cancellations, matching them to the original
reports.
\item Keep or drop each condition explicitly, by purpose: all executed
volume for a volume total; only regular, in-sequence trades for a last
price, a high or a low; auctions separately.
\item Flag prints far from the prevailing quote or from a rolling median;
inspect before deleting: a flash crash is data, a decimal error is not.
\item Record what was removed and why. The rule is part of the result.
\end{enumerate}
\end{method}

\begin{omfigure}
\begin{tikzpicture}
\begin{axis}[width=10.5cm,height=4.8cm,
  xlabel={trade number},ylabel={price},
  tick label style={font=\small},label style={font=\small},
  xmin=0,xmax=16,ymin=0,ymax=115,grid=major,grid style={gray!25},
  scatter/classes={1={mark=*,omDef,mark size=1.8pt},0={mark=x,omThm,mark size=3.5pt,very thick}},
  legend columns=2,legend style={font=\footnotesize,at={(0.5,-0.32)},anchor=north,draw=none,fill=none,/tikz/every even column/.append style={column sep=8pt}},
  table/col sep=comma]
\addplot[scatter,only marks,scatter src=explicit symbolic] table[x=seq,y=price,meta=kept]{figdata/markets-1/28-reading-market-data/badtick.csv};
\legend{kept,removed by the filter}
\end{axis}
\end{tikzpicture}

\omcaption{Fifteen trade reports in one stock. The automatic filter removes
two, the print at 10.01 and a cancelled trade that looks perfectly normal;
the late report at 99.80 passes it. With all fifteen the volume-weighted
average price is 98.40; the true figure is near 100. Data: the weekend
problem.}\label{fig:m1:reading-market-data:badtick}
\end{omfigure}

\section{Timestamps}

\begin{definition}[Exchange and receive timestamps]\label{def:m1:reading-market-data:ts}
The \emph{exchange timestamp}\index{exchange timestamp} is written by the
publisher's matching engine or gateway when the event occurs. The
\emph{receive timestamp}\index{receive timestamp} is written by the
recipient's capture system when the message arrives. The first orders events
within one venue; only the second says what the recipient could have known,
and when.
\end{definition}

A backtest that uses \omterm{def:m1:reading-market-data:ts}{exchange timestamps} assumes the firm learned of each
event at the instant it happened. The difference, tens of microseconds in
the same building and milliseconds between cities, is the whole of some
strategies' edge (One Quant Book 10). Two venues' exchange clocks are not
the same clock: comparing their timestamps to decide which quote moved
first is valid only to the accuracy of their synchronisation.

\begin{dated}{2026-09}{How accurate clocks must be}\label{dat:m1:reading-market-data:clocks}
\emph{European Union.} The regulatory technical standard on clock
synchronisation requires members of a trading venue that use high-frequency
algorithmic techniques to keep their business clocks within 100
microseconds of UTC and to timestamp with a granularity of a microsecond or
better. \emph{United States.} The consolidated audit trail plan requires
exchanges to keep their clocks within 100 microseconds of the national
standards institute's time, and \omterm{def:m1:the-sell-side:sell-side}{broker-dealers} within 50 milliseconds (one
second for clocks used only for manual events).
\end{dated}

\begin{omfigure}
\begin{tikzpicture}
\begin{axis}[width=10cm,height=4.6cm,
  xlabel={mean latency jitter (microseconds)},ylabel={out of order (\%)},
  tick label style={font=\small},label style={font=\small},
  xmin=0,xmax=5200,ymin=0,ymax=36,grid=major,grid style={gray!25},table/col sep=comma]
\addplot[omDef,very thick,mark=*,mark size=1.4pt] table[x=jitter_us,y=inversions_pct]{figdata/markets-1/28-reading-market-data/reorder.csv};
\end{axis}
\end{tikzpicture}

\omcaption{The chapter's sample feed, replayed with a receive time equal to
the exchange time plus 200 microseconds plus an exponential jitter. Sorting
by receive time misorders 1\% of messages at a jitter of 100 microseconds and
a third at 5 milliseconds. Within one feed, order by \omterm{def:m1:reading-market-data:seq}{sequence number}, never
by either clock. Data: the tutorial.}\label{fig:m1:reading-market-data:reorder}
\end{omfigure}

\section{Symbology and reference data}

\begin{definition}[Symbology and reference data]\label{def:m1:reading-market-data:refdata}
\emph{Symbology}\index{symbology} is the set of identifiers under which one
instrument is known: exchange ticker, vendor codes, national and
international security numbers, a feed's numeric code for the day.
\emph{Reference data}\index{reference data} is everything about the
instrument that is not a price: identifiers and their history, listing venue,
currency, tick size, lot size, multiplier, expiry, corporate actions,
trading calendar.
\end{definition}

A ticker is not an identifier. Tickers are reused after delistings, changed
on mergers, and differ by venue and by vendor; a feed's locate code is valid
for one day. Research tables must be keyed by a permanent internal
identifier, with every external symbol mapped to it by date range. The
commonest silent error in equity research, after look-ahead, is a
time series that is two different companies joined at the date one took the
other's ticker. The builds of chapters~\ref{ch:m1:shares-corporate-actions-indices},
\ref{ch:m1:futures-contracts-and-exchanges} and
\ref{ch:m1:option-contracts-and-exchanges} are this chapter's \omterm{def:m1:reading-market-data:refdata}{reference data}.

\section{From raw feed to research table}

\begin{definition}[Normalisation]\label{def:m1:reading-market-data:norm}
\emph{Normalisation}\index{normalisation} is the conversion of each venue's
native messages into one internal representation: common message types,
integer prices in a common unit, a common clock field, permanent instrument
identifiers, explicit condition flags.
\end{definition}

\begin{omfigure}
\begin{tikzpicture}
\begin{axis}[width=11cm,height=5cm,
  xlabel={seconds since the first message},ylabel={price},
  tick label style={font=\small},label style={font=\small},
  ymin=99.945,ymax=100.065,grid=major,grid style={gray!25},
  yticklabel style={/pgf/number format/fixed,/pgf/number format/fixed zerofill,/pgf/number format/precision=2},
  legend columns=3,legend style={font=\footnotesize,at={(0.5,-0.3)},anchor=north,draw=none,fill=none,/tikz/every even column/.append style={column sep=8pt}},
  table/col sep=comma]
\addplot[omDef,thick,const plot] table[x=seconds,y=bid]{figdata/markets-1/28-reading-market-data/l1.csv};
\addplot[omThm,thick,const plot] table[x=seconds,y=ask]{figdata/markets-1/28-reading-market-data/l1.csv};
\addplot[only marks,mark=*,mark size=1.2pt,black] table[x=seconds,y=price]{figdata/markets-1/28-reading-market-data/tape.csv};
\legend{best bid,best ask,trades}
\end{axis}
\end{tikzpicture}

\omcaption{\omterm{def:m1:reading-market-data:levels}{Level 1} and the tape, rebuilt from 640 messages of the chapter's
binary sample by the build's decoder and book. Every execution of a displayed
order prints at the bid or at the ask, which is how its aggressor is known;
the prints between them are trades against hidden orders. Data: the
chapter's build.}\label{fig:m1:reading-market-data:l1}
\end{omfigure}

The pipeline has four stages, and research should know which one its table
came from. \emph{Capture}: every packet, with a hardware \omterm{def:m1:reading-market-data:ts}{receive timestamp},
written unmodified. \emph{Decode and normalise}: the build below.
\emph{Book building}: apply messages in sequence, detect gaps, emit \omterm{def:m1:reading-market-data:levels}{level 2}
and \omterm{def:m1:reading-market-data:levels}{level 1} as derived streams. \emph{Tables}: trades with conditions and
aggressor side; quotes; bars, which are a lossy summary chosen for a
purpose. Each stage must be replayable from the one before, bit for bit:
when a result looks wrong, the question ``is it the data?'' must have an
answer.

\section{Tutorial: a binary feed into a book and a tape}\label{tut:m1:reading-market-data:decode}

\begin{tutorial}
\textbf{Goal.} Decode a framed binary order-by-order feed, maintain the book,
produce \omterm{def:m1:reading-market-data:levels}{level 1} and a tape with aggressor side, and measure what receive-time
jitter does to message order. \textbf{End state:} the three data figures of
this chapter, from \texttt{code/firm/feed/data/sample.itch} (2\,020 messages,
65\,842 bytes).
\begin{tutsteps}
\item \textbf{Framing and header.} Length prefix, type, then the fields every
message shares. A length that does not fit the type is an error, not a
hint.
\omcode{code/firm/feed/firm_feed.py}{58}{74}{Python: frame, check, read the common header.}\label{lst:m1:reading-market-data:py}
\item \textbf{The same in C++}, over a span of bytes, without copies.
\omcode{code/firm/feed/cpp/firm_feed.hpp}{47}{59}{C++: the decode loop's framing and header.}\label{lst:m1:reading-market-data:cpp}
\item \textbf{And in Rust}, as an iterator that yields a result per message
and stops after the first error.
\omcode{code/firm/feed/rust/src/lib.rs}{53}{79}{Rust: one step of the decoder.}\label{lst:m1:reading-market-data:rs}
\item \textbf{The book.} An execution carries no price: it is looked up.
\omcode{code/firm/feed/firm_feed.py}{146}{168}{Applying one message to the book and the tape.}\label{lst:m1:reading-market-data:apply}
\end{tutsteps}
\textbf{What to change next.} Drop one add-order message from the sample and
watch the errors the book records, and how long it takes before an execution
refers to the missing order. Then add a replace message (cancel and add
with a new reference, losing priority) to all three decoders.
\end{tutorial}

\section{Build: the feed normaliser}\label{bld:m1:reading-market-data:feed}

\begin{build}
\bfield{Purpose} The entry point of all market data into the miniature
firm: the simulated exchange of One Quant Book 10 publishes this format, the
strategies of Book 11 consume the book, and the research tables of Book 7 are
written from the tape.
\bfield{Interface} Python: \texttt{encode(msg)}, \texttt{decode(buffer)},
\texttt{Book.apply(msg)}, \texttt{Book.best(locate)}, \texttt{Book.depth(locate,
side, n)}, \texttt{Book.trades}, \texttt{Book.errors}, \texttt{summary(book,
locate)}. C++20 (header only): \texttt{firm::feed::decode(span, callback)},
\texttt{Book}, \texttt{Summary}. Rust: \texttt{decode(\&[u8])} returning an
iterator of results, \texttt{Book}, \texttt{Summary}.
\bfield{Rules} Big-endian integers, prices in ten-thousandths, 48-bit
nanosecond timestamps, two-byte length framing. A truncated frame or a
length inconsistent with the type stops decoding with an error. The book
never guesses: an unknown order reference, an over-execution, a duplicate
add or a timestamp that goes backwards is recorded and counted. The three
implementations must produce the same summary on the shared sample.
\bfield{Acceptance tests} \texttt{code/firm/feed/tests/} (round trip, garbage,
a book by hand, recorded errors, reproducibility of the sample);
\texttt{cpp/firm\_feed\_test.cpp} and the Rust crate's tests (the shared
sample's summary, truncation, a bad length, an unknown order).
\bfield{Stretch} \omterm{def:m1:reading-market-data:seq}{Sequence numbers} with gap detection and A/B line
arbitration; a snapshot message for recovery; a benchmark of messages per
second in each language (One Quant Book 13 starts from there).
\end{build}

\begin{omsources}
\item Nasdaq, \emph{Nasdaq TotalView-ITCH 5.0} specification (revision of 28
April 2023).
\item Commission Delegated Regulation (EU) 2017/574 (clock synchronisation,
``RTS 25'').
\item CAT NMS Plan, frequently asked questions on clock synchronisation
(R1); FINRA Rule 6820.
\end{omsources}

\section{Exercises}

\begin{exercise}[$\star$]\label{exo:m1:reading-market-data:1}
From the \omterm{def:m1:reading-market-data:levels}{level 3} view of \cref{fig:m1:reading-market-data:levels} an execution
of 300 shares arrives for order 1001, then a cancel of 100 shares for order
1002. Give levels 2 and 1 afterwards. At what price did the trade print and
who was the aggressor?
\end{exercise}

\begin{exercise}[$\star$]\label{exo:m1:reading-market-data:2}
A price field contains the four bytes \texttt{00 0F 41 DC} (hexadecimal),
big-endian, with four implied decimals. Give the price. What would a
little-endian reader make of it?
\end{exercise}

\begin{exercise}[$\star$]\label{exo:m1:reading-market-data:3}
The chapter's sample has 1\,102 adds, 311 executions, 114 cancels, 473 deletes
and 20 hidden trades. Verify the file's size from the message lengths and
the framing.
\end{exercise}

\begin{exercise}[$\star\star$]\label{exo:m1:reading-market-data:4}
A timestamp of six bytes counts nanoseconds since midnight. What is the
largest time of day it can hold? Why is four bytes not enough, and what is
gained over eight?
\end{exercise}

\begin{exercise}[$\star\star$]\label{exo:m1:reading-market-data:5}
Venue A timestamps a quote change at 10:00:00.000\,150 and venue B at
10:00:00.000\,190. Both venues' clocks are within 100 microseconds of UTC.
Can you say which moved first? What would you need?
\end{exercise}

\begin{exercise}[$\star\star$]\label{exo:m1:reading-market-data:6}
A receiver sees \omterm{def:m1:reading-market-data:seq}{sequence numbers} 5\,001, 5\,002, 5\,004 on line A and 5\,001,
5\,003, 5\,004 on line B. What does it do? Now both lines skip 5\,003: what
must happen before the book is used again, and what may the strategy do
meanwhile?
\end{exercise}

\begin{exercise}[$\star\star\star$]\label{exo:m1:reading-market-data:7}
\emph{Coding.} With \texttt{receive\_times} and \texttt{inversions} on the
sample, reproduce the out-of-order percentages at jitters of 100 and 1\,000
microseconds. Then halve the message rate (keep every second message) and
explain the change.
\end{exercise}

\begin{exercise}[$\star\star\star$]\label{exo:m1:reading-market-data:8}
\emph{Find the flaw.} A researcher downloads one-minute bars for 3\,000
tickers over fifteen years, keyed by ticker, computes each stock's daily low
as the minimum of the bars' lows, and finds that buying stocks that fell more
than 50\% intraday and recovered by the close is hugely profitable. Name
three data defects, from this chapter, that could produce the result.
\end{exercise}

\section{Problem: The Bad Tick}

\begin{problem}[{Weekend problem --- fifteen prints and a benchmark}]\label{pb:m1:reading-market-data:1}
A desk must report the volume-weighted average price of a stock from 09:30:00
to 09:32:00 as the benchmark of a client order. The tape (sequence, time,
price, shares, condition; the conditions are this book's labels):

\begin{center}\footnotesize
\begin{tabular}{rlrrl@{\qquad}rlrrl}
\toprule
1 & 09:30:00 & 100.02 & 12\,000 & open & 9 & 09:30:52 & 100.09 & 300 & cancelled \\
2 & 09:30:04 & 100.05 & 300 & & 10 & 09:31:03 & 100.11 & 600 & \\
3 & 09:30:09 & 100.04 & 50 & \omterm{def:m1:us-equity-market-structure:lots}{odd lot} & 11 & 09:31:10 & 99.80 & 5\,000 & late \\
4 & 09:30:15 & 100.07 & 500 & & 12 & 09:31:18 & 100.12 & 200 & \\
5 & 09:30:21 & 100.06 & 200 & & 13 & 09:31:25 & 100.13 & 700 & \\
6 & 09:30:30 & 10.01 & 400 & & 14 & 09:31:40 & 100.12 & 100 & \\
7 & 09:30:31 & 100.08 & 400 & & 15 & 09:31:55 & 100.15 & 900 & \\
8 & 09:30:40 & 100.10 & 1\,000 & & & & & & \\
\bottomrule
\end{tabular}
\end{center}

\medskip\noindent\textbf{Part I --- The naive number.}
\begin{enumerate}
\item Give total shares and the volume-weighted average price of all fifteen
prints.
\item Give the period's low from all prints.
\item By how many basis points does the naive average differ from 100.05?
\item What would a volatility estimate from these prints do?
\end{enumerate}

\medskip\noindent\textbf{Part II --- Print by print.}
\begin{enumerate}[resume]
\item Print 6: what is it, how would a filter catch it, and how would you
confirm?
\item Print 9: why does an outlier filter not catch it, and what does?
\item Print 11 was executed at 09:24 and reported at 09:31. Should it be in
the benchmark of an order that started at 09:30?
\item Print 3 is an \omterm{def:m1:us-equity-market-structure:lots}{odd lot}. In or out?
\item Print 1 is the opening auction. The client's order was entered at
09:30:02. In or out?
\end{enumerate}

\medskip\noindent\textbf{Part III --- The corrected numbers.}
\begin{enumerate}[resume]
\item Excluding the erroneous, the cancelled and the late prints, give the
number excluded, the shares, and the average price.
\item Excluding the auction as well, give the average price of continuous
trading.
\item Give the corrected low and high.
\item The client's order bought 2\,000 shares at an average of 100.09. Give
its slippage in basis points against each of the two corrected benchmarks.
\item Which benchmark is right?
\end{enumerate}

\medskip\noindent\textbf{Part IV --- Judgement.}
\begin{enumerate}[resume]
\item The exchange later breaks print 6. How does that reach your database,
and what must your loader do?
\item Why should the cleaning rule be fixed before looking at the order's
result?
\item A vendor's ``cleaned'' bars already exclude some conditions. What must
you obtain from the vendor?
\item Where in the four-stage pipeline does each of the defects of Part II
get handled?
\item State the \emph{named result}: the number of prints to exclude and the
corrected volume-weighted average price.
\item In one sentence: what is a trade report?
\end{enumerate}
\end{problem}

\section{Interview questions}

\begin{interviewq}[$\star$ \iqroles{developer, researcher, trader}]\label{iq:m1:reading-market-data:1}
What is the difference between \omterm{def:m1:reading-market-data:levels}{level 1}, \omterm{def:m1:reading-market-data:levels}{level 2} and \omterm{def:m1:reading-market-data:levels}{level 3} market data?
\end{interviewq}

\begin{interviewq}[$\star$ \iqroles{developer}]\label{iq:m1:reading-market-data:2}
How do you build an order book from an order-by-order feed? What data
structures?
\end{interviewq}

\begin{interviewq}[$\star\star$ \iqroles{developer, researcher}]\label{iq:m1:reading-market-data:3}
Your feed handler detects a gap in \omterm{def:m1:reading-market-data:seq}{sequence numbers}. What happens next?
\end{interviewq}

\begin{interviewq}[$\star\star$ \iqroles{researcher, mle}]\label{iq:m1:reading-market-data:4}
Which timestamp do you use in a backtest, and what goes wrong if you choose
the other?
\end{interviewq}

\begin{interviewq}[$\star\star$ \iqroles{researcher, mle}]\label{iq:m1:reading-market-data:5}
You are given a trades file. What do you check before computing anything?
\end{interviewq}

\begin{interviewq}[$\star\star\star$ \iqroles{developer}]\label{iq:m1:reading-market-data:6}
Design the storage of a firm's captured market data so that any research
result can be reproduced five years later.
\end{interviewq}
